\scanpage{167}
\inserted{Supposons \deleted{$S$} $S$ loc. noeth \uline{ou} $\pi$ artinien [condition des chaînes descendantes]}

(ii) Tout élément $c\in\unsure{\Phi|\varphi}(\pi)$ est \unsure{majoré|minoré} par
un élément \unsure{minimal}, $c'\in\Phi(\pi')$.

% slanted
\margin{\unsure{par un} \unsure{ss-gr.} \unsure{et} \illegible{1} \uline{\unsure{est-il|minimal}} \unsure{NB|2e}}

En effet, si $c$ correspond à $P^{\unsure{\pi}},y$,
supposons que l'on puisse \unsure{réduire|relever} à $P',y'$ \inserted{\unsure{$P'$ (groupe $\pi'\subset\pi$)}} \deleted{\unsure{considérant les images connexes de} \unsure{$y$}.}
et $(P'',y'')$ [groupe $\pi''\subset\pi$].

\margin{$\pi/{\pi'\cap\pi''}\to\pi/\pi'\times\pi/\pi''$}
Considérons $P/{\pi'\cap\pi''}\to P/\pi'\times P/\pi''$
c'est une imm. ouverte et fermée, d'autre
part on a une section du $2^{\mathrm{o}}$ \unsure{membre}
qui \deleted{est} au pt $(y',y'')$ est dans l'image
du premier, elle y est donc partout.

Donc on peut restreindre le groupe structural
à $\pi'\cap\pi''$. Donc les groupes auxquels
on peut restreindre forment un ens
filtrant décroissant, $\pi_i,\ \unsure{\ldots}$, donc $P_i\subset P$
\inserted{forment un syst. projectif.}

Si \deleted{\unsure{les morphismes}} $\pi$ est artinien, on a gagné : le \unsure{systême|système} des $\pi_i$ est
stationnaire. Si $S$ est loc. noeth, on
\deleted{\unsure{les morphismes} de transition $P_j\mathrel{-}P_i$ sont
affines. Donc on peut former le
préschéma $\varprojlim_i P_i$, qui est donc affine sur
$P$, \deleted{\unsure{limite,} \unsure{il est}} \inserted{je dis que l'un des \unsure{objets|sous-groupes}} est
ouvert et fermé, et est un fibré}
% archive sheet number: 166
